% Pursuing Stacks (cotes 134-2 to 134-8): archive cote 134-4, site page 68 (PDF page 69), typescript page 156.
% typescript. Draft compared: ps3.tex:3499-3572.
% AI-transcribed by Claude Opus 5.5 (claude-opus-5-5, adaptive thinking, effort high), Message Batches API, batch ps2 (msgbatch_01NVwr1Yuo8qJf5rykeGBXPK), 2026-11-12; not yet checked by eye.
\scanpage{134-4/68}
\mspage{156}
assuming $\underline{W} = \over{W}{\mathbb{W}}_{(Cat)}$ .

It is in the form of (14) that \unsure{weak|wrak} $\underline{W}$-asphericity of (12) will actu\hand{al}-%
ly be checked, whereas it will be \uline{used} just by the fact that (12) is in
$\underline{W}$ .

6.3.

\margin{62} I finally stopped with the notes last night, \xout{xbxI} by the time \inserted{when} I
started feeling a little unconfortable. A few minutes of reflection then
were enough to convince me that definitely I had'nt done quite enough
preliminary scratchwork yet on this ``key result'' business, and embarked
overoptimistically upon a ``mise en équation'' of the situation, with the
pressing expectation that a caracterization of weak equivalences should
come out at the \xout{first}\inserted[typed]{same} time. First thing that became clear, was that
the introduction of two different localizing sets of arrows $\underline{W},\underline{W}'$ in (Cat)
was rather silly alas, nothing at all would come out unless supposing
from the very start $\underline{W} = \underline{W}'$ . Indeed, the crucial step \hand{for} \xout{thexx"key}
getting the ``key result'' on test functors we are out for, goes as fol-%
lows.

As the target map
\[ t\colon \underline{Fl}(C) = \underline{Hom}(\hand{\Delta}_1,C) \hand{\longrightarrow} C \]
in (Cat) is clearly a homotopy\inserted{retraction}\xout{xx}, and $i^*$: (Cat)$\hand{\to}$ A\^{} commutes with
products, we do have a good hold on the condition (14) of last corollary\hand{,}
namely that $i^*(t)$ be $\underline{W}_A$ - aspheric - e.g. it is enough that the
contractor $i^*(\hand{\Delta}_1)$ in A\^{} be $\underline{W}_A$-aspheric over $e_{A\hat{}}$ (for instance, it is \inserted{often}
enough it be 0-connected !). In view of the corollary and the condition
f) on $\underline{W}$ (last page)\hand{,} we thus get a very good hold upon the \xout{conditionx} map
\[ (\hand{*}) \qquad A_{/C} \hand{\longrightarrow} A_{//C} \]
in (Cat) being in $\underline{W}$ , and hence on all three maps in the diagram (11)
(page 153) being in $\underline{W}$ . With this in mind, the key step can be stated
as follows :

\uline{Lemma.} Assume that $\underline{W}$ satisfies \xout{all}\inserted[typed]{the} conditions a) \xout{toxxf)}\inserted[typed]{c)e)f)} (pages 150,154,
155), that A is $\underline{W}$-aspheric over e (i.e. AxC$\hand{\to}$C is in $\underline{W}$ for any C, \xout{wh}
which will be satisfied if A is $\hand{\Delta}_1$-contractible in (Cat)\hand{,} for instance
if A has \inserted{a} final or initial \xout{element)} object), and that the objects i(a)
in (Cat) (for any a in A) are contractible (for the canonical homotopy
structure of (Cat), i.e. $\hand{\Delta}_1$-contractible\deleted{s}, or even only for the \xout{smarx}
wider homotopy structure $h_{\underline{W}}$ based on $\underline{W}$-aspheric homotopy intervals).

\margin{68}
